\title{The Era of 1-bit LLMs: All Large Language Models are in 1.58 Bits}

\begin{abstract}
Emerging studies, exemplified by BitNet~\cite{bitnet}, are ushering in a paradigm shift toward 1-bit Large Language Models (LLMs). This paper presents a novel 1-bit LLM variant, \textbf{\text{BitNet b1.58}}, characterized by its use of ternary weights drawn from the set \{-1, 0, 1\} for every parameter. Remarkably, this design attains perplexity and downstream task performance on par with its full-precision Transformer LLM counterparts (FP16 or BF16) of comparable size and trained on equivalent data, yet it delivers notable improvements in latency, memory efficiency, throughput, and energy usage. Notably, the 1.58-bit LLM introduces a revised scaling law and training framework, guiding the development of future LLMs that are simultaneously high-performing and resource-efficient. Additionally, this advancement lays the groundwork for novel computational paradigms and invites the innovation of hardware architectures tailored for 1-bit LLMs.
\end{abstract}

\section{The Era of 1-bit LLMs}

The AI landscape has recently witnessed a surge in the scale and proficiency of Large Language Models (LLMs). While these advances have enabled impressive outcomes across diverse natural language processing applications, the exponential growth in model size presents obstacles for real-world deployment and amplifies worries regarding energy usage and sustainability. 

A widely adopted strategy to tackle these issues is post-training quantization, which produces low-precision models suited for inference~\cite{smoothquant, gptq, quip, quip_sharp}. By compressing the numerical representation of weights and activations, this method greatly lowers the storage and computational burdens of LLMs. The evolution of quantization methods reflects a movement from 16-bit down to lower bit-widths, notably including 4-bit configurations~\cite{gptq, awq}. Nevertheless, although post-training quantization is prevalent in production systems, it often yields sub-optimal solutions.

Recent advancements in 1-bit model designs, exemplified by BitNet~\cite{bitnet}, have opened up new avenues for lowering the operational expenses of LLMs while preserving their accuracy. Standard LLMs typically utilize 16-bit floating-point representations (either FP16 or BF16), with matrix multiplication dominating the computational workload. The bulk of computation in these architectures stems from floating-point additions and multiplications. In comparison, BitNet restructures matrix multiplication to rely solely on integer additions, dramatically reducing the energy consumption required for large language model computation. Since energy usage often limits the computational performance of hardware accelerators, such improvements in efficiency also yield faster processing times.

Beyond raw computation, the act of moving model weights from external DRAM into local, on-chip accelerator memory (such as SRAM) can present a major inference bottleneck. Although enlarging SRAM capacities has been explored as a way to enhance throughput, this approach incurs considerably greater costs relative to DRAM. The memory footprint—both in terms of capacity and bandwidth—of 1-bit LLMs is substantially smaller than that of full-precision counterparts. This translates into a significant reduction in both the cost and duration required to transfer model parameters, ultimately leading to more expedient and cost-effective inference.

Here, we present a notable 1-bit LLM alternative, \textbf{\text{BitNet b1.58}}, characterized by all parameters being ternary values from the set \{-1, 0, 1\}. This is achieved by augmenting the original BitNet’s binary parameterization with an explicit zero, producing an effective representation of 1.58 bits in binary terms. \text{BitNet b1.58} maintains every advantage of the prior 1-bit BitNet, including its innovative computation strategy that minimizes the use of multiplication in matrix multiplications and lends itself to extensive computational optimizations. Furthermore, the energy usage mirrors that of the original BitNet, and compared to FP16 LLM standards, it provides marked improvements in memory efficiency, processing throughput, and latency. On top of these, \text{BitNet b1.58} introduces two unique benefits. Firstly, with the zero value incorporated into its weight space, it explicitly enables feature filtering, thereby significantly enhancing its expressive capacity and performance relative to traditional 1-bit LLMs. Secondly, empirical results indicate that \text{BitNet b1.58} can attain parity with full-precision (FP16) models in both perplexity and end-task metrics from the 3B parameter scale onwards, given matched configurations (such as model size and amount of training data).

\section{BitNet b1.58}

\text{BitNet b1.58} is constructed upon the BitNet framework, a variant of the Transformer model where the standard \emph{nn.Linear} operation is substituted by \emph{BitLinear}. This model is trained from the ground up, utilizing weights quantized to 1.58 bits and activations to 8 bits. Relative to the original BitNet, several alterations have been introduced, which are outlined below.

\paragraph{Quantization Function.} To ensure the weights are restricted to the set \{-1, 0, +1\}, an \emph{absmean} quantization function is employed. The procedure involves normalizing the weight matrix by its mean absolute value, followed by rounding each element to its closest member in \{-1, 0, +1\}:

\begin{equation}
    \widetilde{W} = \mathrm{RoundClip}(\frac{W}{\gamma + \epsilon}, -1, 1),
\end{equation}

\begin{equation}
    \mathrm{RoundClip}(x, a, b) = \max(a, \min(b, \mathrm{round}(x))),
\end{equation}

\begin{equation}
    \gamma = \frac{1}{nm}\sum_{ij} |W_{ij}|.
\end{equation}

For activation quantization, the methodology remains consistent with the implementation found in BitNet, but with one key distinction: activations are no longer rescaled to $[0, Q_b]$ prior to applying non-linearities. Instead, activations for each token are scaled within $[-Q_b, Q_b]$, effectively removing zero-point quantization. This approach simplifies both implementation and system-level optimization, and in our empirical analysis, results in only negligible changes in performance.

\paragraph{LLaMA-alike Components.}

With LLaMA~\cite{llama, llama2} having established itself as a core architecture within the open-source large language model community, \text{BitNet b1.58} intentionally incorporates LLaMA-inspired modules to foster compatibility. Specifically, it integrates RMSNorm~\cite{rmsnorm}, SwiGLU~\cite{swiglu}, and rotary embedding~\cite{rope}, and excludes bias terms throughout the network. This ensures that \text{BitNet b1.58} can be readily incorporated into widely used open-source platforms (such as Huggingface, vLLM~\cite{vllm}, and llama.cpp\footnote{\url{https://github.com/ggerganov/llama.cpp}}) with little to no modification required.

\section{Results}

We evaluated \text{BitNet b1.58} alongside our re-implemented FP16 \text{LLaMA LLM} across several model sizes. Both models were pre-trained on the RedPajama dataset~\cite{redpajama} with an identical token count of 100 billion, enabling equitable conditions for assessment. Zero-shot evaluations spanned multiple language understanding benchmarks, such as ARC-Easy~\cite{arc}, ARC-Challenge~\cite{arc}, Hellaswag~\cite{hellaswag}, Winogrande~\cite{winoGrande}, PIQA~\cite{piqa}, OpenbookQA~\cite{openbookqa}, and BoolQ~\cite{boolq}. Additionally, we monitored validation perplexity using WikiText2~\cite{wikitext2} and C4~\cite{c4} corpora.

For performance profiling, we measured the GPU memory consumption and latency during inference for both \text{LLaMA LLM} and \text{BitNet b1.58}. Measurements were collected leveraging the FasterTransformer\footnote{\url{https://github.com/NVIDIA/FasterTransformer}} library, known for efficient inference of LLMs on GPUs. In the case of \text{BitNet b1.58}, the evaluation incorporated the Ladder~\cite{ladder} 2-bit kernel. Latency was quantified as the average time per generated token, reflecting the dominant component of inference computation.

Perplexity scores and resource usage for both \text{BitNet b1.58} and \text{LLaMA LLM} are detailed in Table~\ref{tab:ppl}. The findings indicate that beginning at the 3B parameter scale, \text{BitNet b1.58} closely approaches the perplexity achieved by the full-precision \text{LLaMA LLM}, while realizing a 2.71$\times$ increase in speed and a 3.55$\times$ reduction in GPU memory usage. Notably, the 3.9B variant of \text{BitNet b1.58} attains 2.4$\times$ lower latency and 3.32$\times$ less memory overhead compared to \text{LLaMA LLM} 3B, and yields notably superior accuracy.

The comprehensive zero-shot accuracy results for downstream tasks appear in Table~\ref{tab:zero-shot}. Evaluation followed the methodology of \emph{lm-evaluation-harness}\footnote{\url{https://github.com/EleutherAI/lm-evaluation-harness}}. The data demonstrate that as model size increases, the performance differential between \text{BitNet b1.58} and \text{LLaMA LLM} diminishes. Crucially, from the 3B model size, \text{BitNet b1.58} achieves parity with its full-precision counterpart. Consistent with perplexity trends, the end-task outcomes show that \text{BitNet b1.58} at 3.9B not only surpasses \text{LLaMA LLM} 3B but does so with significantly reduced computational cost. Overall, these results position \text{BitNet b1.58} as delivering a Pareto improvement over leading LLM baselines.

\paragraph{Memory and Latency}

We expanded our experiments to encompass model sizes of 7B, 13B, and 70B, analyzing their associated computational costs. As shown in Figure~\ref{fig:latency-memory-energy}, latency and memory consumption both demonstrate clear scaling patterns, with observed acceleration growing in tandem with model size. Notably, at 70B parameters, \text{BitNet b1.58} achieves a 4.1-fold improvement in speed relative to the \text{LLaMA LLM} baseline. This enhancement is attributed to the increasing contribution of the \emph{nn.Linear} operation's time overhead as model parameters rise. Similarly, memory usage reflects this scaling trend—since the embedding layers remain in full precision, their proportion of total memory diminishes as the model increases in size. All measurements for latency and memory employed a 2-bit kernel, indicating further optimizations may be possible to minimize costs.

\paragraph{Energy}

We proceeded to estimate the energy expenditure for arithmetic operations within both \text{BitNet b1.58} and \text{LLaMA LLM}, concentrating on matrix multiplication given its predominance in large language model computations. Figure~\ref{fig:energy} details the respective contributions to energy usage: \text{BitNet b1.58} primarily employs INT8 addition, whereas \text{LLaMA LLM} relies on both FP16 addition and FP16 multiplication. Based on the energy model referenced in~\cite{energycost, pokebnn}, matrix multiplication using \text{BitNet b1.58} on 7nm chips is estimated to require 71.4 times less energy for arithmetic operations than the baseline. Additionally, we report the total energy usage for full model inference with 512 tokens. The data indicate that, as model size increases, the relative energy efficiency advantage of \text{BitNet b1.58} over the FP16-based \text{LLaMA LLM} expands. This effect is explained by the growing proportion of computational cost due to \emph{nn.Linear} in larger models, while the contribution from other model components becomes comparatively less significant.

\paragraph{Throughput}

To evaluate throughput, we compared \text{BitNet b1.58} and \text{LLaMA LLM} models, each with 70B parameters, across two 80GB A100 GPUs using pipeline parallelism~\cite{gpipe}, which allows \text{LLaMA LLM} 70B to be executed within available hardware constraints. By incrementally increasing the batch size up to the maximum supported by GPU memory (with sequence length fixed at 512), we observed in Table~\ref{tab:throughput} that \text{BitNet b1.58} 70B is capable of accommodating a batch size up to 11 times greater than that of \text{LLaMA LLM}, ultimately delivering an 8.9-fold improvement in throughput.

\textbf{\text{BitNet b1.58} introduces a novel scaling law relating model effectiveness to inference efficiency.} Based on data presented in Figure~\ref{fig:latency-memory-energy} and \ref{fig:energy}, the following correspondences emerge for 1.58-bit and 16-bit model configurations:

\begin{itemize}
    \item In terms of latency, memory consumption, and energy use, 13B BitNet b1.58 surpasses 3B FP16 LLM in efficiency.
    \item For 30B BitNet b1.58, it outperforms 7B FP16 LLM across latency, memory footprint, and energy metrics.
    \item 70B BitNet b1.58 demonstrates greater efficiency in all three measures compared to 13B FP16 LLM.
\end{itemize}

\paragraph{Training with 2T Tokens}

Training token count is an essential component in LLM scaling. To assess the token scalability of \text{BitNet b1.58}, we pretrained a \text{BitNet b1.58} model using 2T tokens, adhering to the StableLM-3B~\cite{StableLM-3B-4E1T} data pipeline, representative of the best-performing 3B open-source model. Evaluations were conducted using a comprehensive benchmark suite, including Winogrande~\cite{winoGrande}, PIQA~\cite{piqa}, SciQ~\cite{sciq}, LAMBADA~\cite{lambada}, and ARC-easy~\cite{arc}. Zero-shot accuracy results are listed in Table~\ref{tab:2t}, averaging accuracy metrics where both raw and normalized values are reported. StableLM 3b performance at 2T tokens is quoted from its official report. Our empirical analysis demonstrates that \text{BitNet b1.58} consistently delivers superior outcomes across all tasks, highlighting robust generalization for 1.58-bit LLMs.

\section{Discussion and Future Work}

\textbf{1-bit Mixture-of-Experts (MoE) LLMs}

Mixture-of-Experts (MoE) architectures are recognized as an efficient solution for reducing computational cost in LLMs. However, despite achieving lower FLOPs, their large memory requirements and substantial inter-chip communication overhead remain key obstacles for widespread use. The advent of 1.58-bit LLMs addresses these limitations: notably, the smaller memory footprint enables MoE models to operate with fewer hardware devices. In addition, the diminished activation size leads to a significant reduction in network transfer overhead, with the ultimate possibility of eliminating such costs if the full model fits onto a single device.

\textbf{Native Support of Long Sequence in LLMs}

Handling long input sequences has emerged as a fundamental requirement for modern LLMs. A primary bottleneck for inference over extended contexts is the extensive memory needed to maintain KV caches. \text{BitNet b1.58} provides substantial progress towards efficient support for long sequences by shrinking activation sizes from 16 bits to 8 bits, thus enabling the context window to double under fixed computational budgets. Furthermore, this can potentially be reduced further—to 4 bits or less—in future iterations of 1.58-bit LLMs, representing an open avenue for further investigation.

\textbf{LLMs on Edge and Mobile}

Deploying 1.58-bit LLMs may profoundly enhance language model performance on edge and mobile devices, where memory and processing power are at a premium. Such hardware constraints often limit the feasibility and scale of LLM applications in these environments. By reducing memory and energy demands, 1.58-bit LLMs make it possible to implement more advanced functionalities on such platforms, broadening their application landscape. In particular, these models align well with CPUs, the predominant processing units in edge and mobile scenarios, making it feasible for \text{BitNet b1.58} to run efficiently and thereby unlocking greater performance and new use cases for these devices.

\textbf{New Hardware for 1-bit LLMs}

Projects such as~Groq\footnote{\url{https://groq.com/}} illustrate the significant advancements and opportunities in designing specialized hardware (such as LPUs) for LLM acceleration. Building on these trends, we propose and encourage research focused on architecting novel hardware and system designs explicitly tailored for the unique requirements of 1-bit LLMs, to fully leverage the computational innovations introduced by BitNet~\cite{bitnet}.